% results_refinement.tex
\subsection{Response-Free Refinement}

Three functions refine a scale before any respondent is recruited.
\texttt{sfa\_redundancy()} detects locally dependent item pairs with
Unique Variable Analysis---wTO on an EBICglasso
network \citep{christensen-etal-2023}:

{\scriptsize
\begin{verbatim}
> red_uva <- sfa_redundancy(fit)                     # UVA / wTO (EGAnet)
> red_uva
Redundant-item detection
  Method: Christensen et al. (2023)
  Measure: wto | threshold: 0.25
  Redundant pairs: 11 | redundant clusters: 10

  Top redundant pairs:
    A24        ~ A29         overlap=0.462
    N17        ~ N18         overlap=0.402
    O45        ~ O50         overlap=0.374
    N16        ~ N19         overlap=0.361
    E2         ~ E6          overlap=0.339
    A25        ~ A27         overlap=0.331
    A21        ~ A27         overlap=0.320
    C32        ~ C36         overlap=0.309
    C33        ~ C40         overlap=0.286
    O41        ~ O48         overlap=0.268

  Suggested removals (keep one per cluster): E6, N19, N18, A25, A27, A29,
    C36, C40, O41, O44, O45
\end{verbatim}
}

Every displayed pair among the \valUvaPairsWord{} (in \valUvaClustersWord{} clusters) joins items from
one domain; the simpler cosine criterion (threshold 0.85) flagged a
largely overlapping set of \valCosPairsWord{} pairs in \valCosClustersWord{} clusters, led by
N16~$\sim$~N19 (\valCosTopA) and A24~$\sim$~A29 (\valCosTopB)---near-paraphrases, the
``cheating by repeating'' that inflates internal consistency without
broadening coverage \citep{christensen-etal-2023}.

\texttt{sfa\_simplify()} builds candidate short forms from the items
nearest each construct anchor \citep{jung-seo-2025, wang-etal-2026}:

{\scriptsize
\begin{verbatim}
> short <- sfa_simplify(fit, target_n = 5, method = "anchor")   # 50 -> 25 items
> short
Scale simplification (response-free)
  Method: centroid/medoid item selection (in the spirit of Wang et al. 2026;
    Jung & Seo 2025)
  Note: candidate short form -- validate psychometrically before use
  Selection: anchor | groups: theoretical | target_n = 5 per group
  Items: 50 -> 25
  Factors retained (parallel analysis): 5 -> 3
  Structure recovery vs theory (NMI / ARI):
    full:    NMI=0.527  ARI=0.402
    reduced: NMI=0.810  ARI=0.571
  Items kept per construct:
    Extraversion         5
    Neuroticism          5
    Agreeableness        5
    Conscientiousness    5
    Openness             5

  Dropped 25 item(s); see $drop. Reduced fit in $reduced_fit.
> short$keep
 [1] "E2"  "E3"  "E5"  "E8"  "E10" "N11" "N15" "N16" "N17" "N19" "A21" "A22"
[13] "A24" "A28" "A29" "C33" "C34" "C35" "C37" "C40" "O42" "O43" "O45" "O48"
[25] "O50"
\end{verbatim}
}

Halving the scale to 25 items \emph{improved} theoretical-partition
recovery (NMI \valSfNmiFull{} to \valSfNmiRed; ARI \valSfAriFull{} to \valSfAriRed): pruning semantically
peripheral items sharpens boundaries the full pool blurs. Parallel
analysis now suggested \valSfKRedWord{} factors; per the printed caveat, a semantic
short form is a candidate for psychometric validation, not a finished
instrument.

\texttt{sfa\_item\_fit()} vets newly written items against an existing
fit, extending \citet{wulff-mata-2025}. Live
embedding confines the demonstration to the default-model (0.6B) space;
the reference fit omits the scoring column for the \emph{anti-topic}
reasons given in the Method section.

{\scriptsize
\begin{verbatim}
> fit06 <- sfa(big5_df[, c("code", "item", "factor")],
+              embeddings = ref06$embeddings,
+              encoding = "mean_centered_pearson", nfactors = 5)
> cand <- c("I make friends easily.",
+           "I am the life of every party.",
+           "I rarely feel anxious or depressed.")
> ifit <- sfa_item_fit(fit06, cand, model = "Qwen/Qwen3-Embedding-0.6B",
+                      redundancy_cutoff = 0.85)
> ifit
Candidate item fit
  Method: Wulff & Mata (2025), extended

Candidate 1: "I make friends easily."
  [similarity-by-construct table omitted]
  Best construct : Extraversion
  Cross-loading  : gap to 2nd (Neuroticism) = 0.01  ->  ambiguous -
    cross-loads with Neuroticism
  Strength       : about typical for Extraversion (items 0.74 vs avg 0.75)
  Redundancy     : closest item E3 "I feel comfortable around people."
    (0.78) -> distinct
  Verdict        : cross-loads (Extraversion vs Neuroticism) - poor
    discrimination

Candidate 2: "I am the life of every party."
  [similarity-by-construct table omitted]
  Best construct : Extraversion
  Redundancy     : closest item E1 "I am the life of the party." (0.90)
    -> near-duplicate
  Verdict        : redundant (near-duplicate of an existing item)

Candidate 3: "I rarely feel anxious or depressed."
  [similarity-by-construct table omitted]
  Best construct : Neuroticism
  Strength       : weaker than average for Neuroticism (items 0.75 vs
    avg 0.81)
  Verdict        : weak match - fits Neuroticism worse than its existing
    items (may not belong to this scale)
\end{verbatim}
}

All three planted candidates drew the intended verdicts: cross-loading
(Extraversion vs.\ Neuroticism gap \valIfGap), near-duplication of E1 (cosine
\valIfNearest, above the .85 cutoff), and a weaker match to its best construct
than that construct's existing items.
